\documentclass[11pt]{article}

\usepackage[margin=1in]{geometry}
\usepackage{microtype}
\usepackage[hidelinks]{hyperref}
\setlength{\parindent}{0pt}
\setlength{\parskip}{0.5em}

\hypersetup{
  pdftitle={Cover Letter to Physical Review D},
  pdfauthor={Ronald Bibb},
  pdfsubject={Submission of Normal-Form Resolution of a Free-Boundary MOTS Saddle-Node in Dynamical Five-Dimensional Gravity}
}

\begin{document}

\begin{flushright}
Ronald Bibb\\
Independent Researcher\\
Lilburn, Georgia, USA\\
\href{mailto:ronbibb@gmail.com}{ronbibb@gmail.com}\\[0.75em]
August 24, 2026
\end{flushright}

Editors\\
\textit{Physical Review D}

Dear Editors:

I submit the manuscript ``Normal-Form Resolution of a Free-Boundary MOTS
Saddle-Node in Dynamical Five-Dimensional Gravity'' for consideration as a
Research Article in \textit{Physical Review D}.

The manuscript addresses a boundary-sensitive question in black-hole horizon
theory and numerical relativity.  Marginally outer trapped surface (MOTS) pair
creation is already understood as a saddle-node mechanism for closed surfaces.
Here the surface terminates on a physical brane, so the contact condition
changes the domain of the stability operator and the corresponding discrete
adjoint projection.  The ambient five-dimensional geometry is also generated
by a nonlinear Einstein--scalar initial-boundary-value evolution rather than
prescribed as an exact metric family.  Existing closed-MOTS examples and
free-boundary stability or rigidity theory therefore do not by themselves
establish that the fold hypotheses are realized in this operator problem.

The principal result is a spatially and temporally closed numerical resolution
of that critical structure in the $SO(3)$-symmetric sector.  An independently
constructed parent sequence evolves from no accepted brane-capped MOTS to two
opposite-stability branches on three successive spacetime grids.
Pseudo-arclength continuation connects the branches through a common critical
surface.  The principal stability eigenvalue is bracketed through zero with an
isolated sign-definite mode; endpoint-eliminated discrete adjoint projections
give nonzero transversality and quadratic coefficients; and the invariant
proper-area separation follows the expected square-root law.  An independent
G10 half-timestep evolution repeats the parent and continuation and yields an
overlapping critical-time bracket with a free exponent near one half.

The contribution is not a claim to a new universal bifurcation mechanism.  It
is a controlled numerical realization of the normal form in a dynamical
free-boundary MOTS problem whose physical Robin condition, non-self-adjoint
operator, and evolved higher-dimensional geometry materially alter the
calculation.  This intersection of general relativity, horizon theory,
extra-dimensional gravity, and computational physics falls directly within
the scope of \textit{Physical Review D} and should be of interest to readers
working on black-hole dynamics and numerical relativity.

The manuscript carefully limits the inference to finite-resolution,
quasi-local evidence: it does not claim a continuum theorem, event-horizon
reconstruction, nonsymmetric nonprincipal spectrum, or phase diagram.  The
supporting numerical records, frozen protocols, source manifests, and tests
are publicly archived at
\href{https://doi.org/10.5281/zenodo.22087164}{doi:10.5281/zenodo.22087164},
with version-controlled source at
\href{https://github.com/RonBibb/free-boundary-mots-saddle-node}{GitHub}.
The manuscript is submitted as a regular Research Article and is not part of
a joint submission.

Thank you for considering this work.

Sincerely,

Ronald Bibb

\end{document}
